\documentclass[11pt]{article}

\usepackage[margin=1in]{geometry}
\usepackage{amsmath,amssymb}
\usepackage{empheq}
\usepackage{bm}
\usepackage{graphicx}
\usepackage{hyperref}
\usepackage{xcolor}

\newcommand{\chieff}{\chi_{\mathrm{eff}}}
\newcommand{\Mc}{\mathcal{M}}
\newcommand{\Mtot}{M_{\mathrm{tot}}}
\newcommand{\Fisher}{\bm{F}}
\newcommand{\Cov}{\bm{C}}
\newcommand{\vth}{\bm{\theta}}
\newcommand{\diag}{\mathrm{diag}}
\newcommand{\sgn}{\mathrm{sgn}}

\title{Optimally Measured Intrinsic Parameters of Gravitational-Wave Sources:\\
Generalisations of $\chi_{\mathrm{eff}}$ and Chirp Mass}
\author{}
\date{\today}

\begin{document}
\maketitle

\begin{abstract}
We determine the linear combinations of a compact-binary's intrinsic parameters
that are \emph{best measured} by a gravitational-wave detector network, as
smooth functions of the true source parameters. Working in the Fisher
information approximation, we marginalise over all other parameters, whiten the
resulting covariance so that parameters of different dimension are comparable,
and diagonalise. The eigenvector of smallest variance defines the best-measured
combination. For aligned spins this combination reproduces the effective spin
$\chieff$ to leading order and we provide a closed-form generalisation. For the
masses, expressed in $(\log m_1,\log m_2)$, the best-measured combination
smoothly interpolates between the chirp mass at low total mass and the total
mass at high total mass.
\end{abstract}

\section{Setup and conventions}

We consider a quasi-circular, aligned-spin compact binary. The frequency-domain
waveform $h(f;\vth)$ depends on the parameter vector
\begin{equation}
  \vth = \bigl(\Mc,\,\eta,\,d_L,\,\theta,\,\phi,\,\iota,\,\psi,\,t_c,\,\phi_c,\,
  \chi_{1z},\,\chi_{2z}\bigr),
\end{equation}
with detector-frame chirp mass $\Mc$, symmetric mass ratio
$\eta=m_1 m_2/(m_1+m_2)^2$, luminosity distance $d_L$, sky position
$(\theta,\phi)$, inclination $\iota$, polarisation $\psi$, coalescence time and
phase $(t_c,\phi_c)$, and the dimensionless aligned spin components
$\chi_{1z},\chi_{2z}$. Source-frame masses are redshifted to the detector frame,
$m_i^{\rm det}=(1+z)\,m_i^{\rm src}$, and the redshift fixes $d_L$ through a
fiducial cosmology. The masses can equivalently be parametrised by
$(m_1,m_2)$ or $(\log m_1,\log m_2)$; we use the last of these for the mass
study (Sec.~\ref{sec:mass}).

\paragraph{Fisher matrix.}
For a network of detectors with one-sided noise PSDs $S_n^{(d)}(f)$, the Fisher
information matrix is
\begin{equation}
  F_{ij} = \sum_d \left( \frac{\partial h}{\partial \theta_i}\,\Big|\,
  \frac{\partial h}{\partial \theta_j} \right)_d,
  \qquad
  (a\,|\,b)_d = 4\,\mathrm{Re}\!\int_{f_{\min}}^{f_{\max}}
  \frac{a^*(f)\,b(f)}{S_n^{(d)}(f)}\,df .
  \label{eq:fisher}
\end{equation}
The waveform derivatives are evaluated by automatic differentiation
(\texttt{gwfast}/JAX), giving smooth, low-noise gradients --- essential for the
eigenvectors to vary smoothly with $\vth$. In the large-SNR (Gaussian-linear)
limit the parameter covariance is $\Cov=\Fisher^{-1}$.

\section{Best-measured directions}
\label{sec:method}

Let $\vth_S \subset \vth$ be an intrinsic sub-space of interest (the spins, or
the masses). Marginalising over all other parameters under the Gaussian
approximation amounts to taking the corresponding sub-block of the full
covariance,
\begin{equation}
  \Cov_S = \bigl[\Fisher^{-1}\bigr]_{S}.
  \label{eq:marg}
\end{equation}
A linear combination $p = \bm{a}^{\!\top}\vth_S$ has variance
$\sigma_p^2 = \bm{a}^{\!\top}\Cov_S\,\bm{a}$. The best-measured combination
minimises this variance, so it is the eigenvector of $\Cov_S$ with the
\emph{smallest} eigenvalue.

\paragraph{Whitening.}
Because the parameters carry different units and natural scales (e.g.\ $\Mc$ vs.\
$\eta$), the raw eigenproblem is scale-dependent. We whiten by the marginal
standard deviations $s_i=\sqrt{(\Cov_S)_{ii}}$, i.e.\ work in the dimensionless
variables $\bm{u}=\bm{D}^{-1}\vth_S$ with $\bm{D}=\diag(s_i)$. The whitened
covariance is the correlation matrix
\begin{equation}
  \Cov_S^{\,w} = \bm{D}^{-1}\Cov_S\,\bm{D}^{-1},
\end{equation}
which we symmetrise and diagonalise,
$\Cov_S^{\,w}\,\bm{b}_j = \lambda_j\,\bm{b}_j$, with
$\lambda_1\le\lambda_2\le\dots$ ordered ascending. The whitened eigenvectors
$\bm{b}_j$ are orthonormal. The corresponding combination in the original
parameters has coefficients
\begin{equation}
  \bm{a}_j = \bm{D}^{-1}\bm{b}_j \;=\; \bm{b}_j / \bm{s},
  \label{eq:phys}
\end{equation}
which we normalise to unit Euclidean norm. The leading vector $\bm{a}_1$ is the
best-measured combination; $\bm{a}_2$ is the next-best, and so on.

\paragraph{Numerical reductions.}
Three equivalent-in-principle reductions are implemented to assess numerical
stability: (i) \emph{marginalise} --- invert the full Fisher, take the sub-block
[Eq.~\eqref{eq:marg}]; (ii) \emph{whiten-then-marginalise} --- whiten the full
Fisher by its diagonal before inversion; (iii) \emph{condition} --- fix all
other parameters and invert the Fisher sub-block. We adopt (i) by default.

\paragraph{Smoothness across the grid.}
The sign and ordering of eigenvectors are arbitrary at each point. We track each
mode along the parameter grid by matching the orthonormal whitened eigenvectors
$\bm{b}_j$ between adjacent points (maximising overlap and fixing the sign for
positive overlap), which is robust through eigenvalue crossings. Matching the
non-orthonormal physical coefficients $\bm{a}_j$ would mis-assign near
crossings.

\section{Spin: generalisation of $\chieff$}
\label{sec:spin}

The effective aligned spin is
\begin{equation}
  \chieff = \frac{\chi_{1z}+q\,\chi_{2z}}{1+q},
  \qquad q \equiv \frac{m_2}{m_1}\le 1,
  \label{eq:chieff}
\end{equation}
i.e.\ the combination of $(\chi_{1z},\chi_{2z})$ with weight ratio
$c_2/c_1 = q$. We compute the best-measured spin direction in the
$(\chi_{1z},\chi_{2z})$ sub-space by the procedure of Sec.~\ref{sec:method} and
report its weight ratio $c_2/c_1$.

Numerically the best-measured direction tracks $\chieff$ closely but lies
systematically slightly below it. A two-parameter fit (single LIGO at design
sensitivity, \texttt{IMRPhenomD}; weak dependence on the other parameters) gives
\begin{empheq}[box=\fbox]{equation}
  \left(\frac{c_2}{c_1}\right)_{\rm best}
  = q\Bigl[\,1 - 0.337\,(1-q) + 0.359\,(1-q)^2\,\Bigr],
  \label{eq:spinfit}
\end{empheq}
with RMS residual $\approx 0.018$ and a maximum deviation from the $\chieff$
prediction $c_2/c_1=q$ of $\approx 0.10$ near $q\sim0.5$. Equation
\eqref{eq:spinfit} reduces to $\chieff$ at equal mass ($q\to1$) and stays within
$\sim10\%$ of it across the range, quantifying the (modest) sense in which
$\chieff$ is the optimally measured aligned-spin parameter.

\section{Mass: chirp mass to total mass}
\label{sec:mass}

We study the masses in the basis $(\log m_1,\log m_2)$, in which the
best-measured combination $w_1\log m_1 + w_2\log m_2$ is a power-law mass
parameter $\propto m_1^{w_1} m_2^{w_2}$ and is scale-free. The Fisher matrix in
this basis is obtained from the $(m_1,m_2)$ Fisher by the diagonal Jacobian
$\Fisher_{\log} = \bm{J}\Fisher\bm{J}$ with $\bm{J}=\diag(m_1,m_2,1,\dots)$.

Two reference combinations bracket the expected behaviour. The chirp mass
\begin{equation}
  \Mc = \frac{(m_1 m_2)^{3/5}}{(m_1+m_2)^{1/5}},
  \qquad
  \log\Mc = \tfrac{3}{5}\bigl(\log m_1+\log m_2\bigr)
            - \tfrac{1}{5}\log(m_1+m_2),
\end{equation}
has gradient weights in $(\log m_1,\log m_2)$ with ratio
\begin{equation}
  \left(\frac{w_2}{w_1}\right)_{\!\Mc}
  = \frac{\partial\log\Mc/\partial\log m_2}{\partial\log\Mc/\partial\log m_1}
  = \frac{3+2q}{2+3q},
  \label{eq:rchirp}
\end{equation}
while the total mass $\Mtot=m_1+m_2$ has
\begin{equation}
  \left(\frac{w_2}{w_1}\right)_{\!\Mtot}
  = \frac{m_2/\Mtot}{m_1/\Mtot} = q .
  \label{eq:rtot}
\end{equation}
At low total mass, where many inspiral cycles lie in band, the phase is
dominated by $\Mc$ and the best-measured combination matches
Eq.~\eqref{eq:rchirp}. At high total mass, where the signal is dominated by the
merger--ringdown ($f\sim 1/\Mtot$), it approaches Eq.~\eqref{eq:rtot}. We model
the crossover as
\begin{empheq}[box=\fbox]{equation}
  \left(\frac{w_2}{w_1}\right)_{\rm best}
  = \bigl[1-s(\Mtot)\bigr]\,\frac{3+2q}{2+3q} + s(\Mtot)\,q,
  \qquad
  s(\Mtot)=\frac{1}{1+\bigl(\Mtot/\Mtot^{\star}\bigr)^{-k}},
  \label{eq:massfit}
\end{empheq}
a sigmoid in $\log_{10}\Mtot$ with transition mass $\Mtot^{\star}$ and sharpness
$k$. The single-detector design fit gives $\Mtot^{\star}\sim 1.5$--$2.5\times
10^{2}\,M_\odot$ for the transition.

\paragraph{Caveats.}
Equation~\eqref{eq:massfit} captures the chirp-to-total trend but is not yet a
high-fidelity fit. The data show (i) a non-monotonic \emph{overshoot} near
$\Mtot\sim100$--$130\,M_\odot$, where $w_2/w_1$ rises above the chirp value
before falling toward the total-mass value --- a real feature the simple sigmoid
cannot represent; and (ii) ill-conditioning at equal mass ($q\to1$), where the
two $(\log m_1,\log m_2)$ eigenvalues become near-degenerate and the
best-measured mode can spuriously flip to the antisymmetric direction
$w_2/w_1\to-1$ (the poorly measured mass-ratio direction). A richer functional
form and explicit regularisation of the $q=1$ degeneracy are required for a
quantitative mass formula; the qualitative chirp$\to$total transition is robust.

\section{Summary}

The best-measured intrinsic parameters, defined as the smallest-variance
eigenvectors of the marginalised, whitened Fisher covariance, generalise the
familiar phenomenological combinations:
\begin{itemize}
  \item \textbf{Spin:} a smooth, $\sim10\%$ correction to $\chieff$,
  Eq.~\eqref{eq:spinfit}, reducing to $\chieff$ at equal mass.
  \item \textbf{Mass:} a power-law mass parameter that interpolates between
  chirp mass [Eq.~\eqref{eq:rchirp}] and total mass [Eq.~\eqref{eq:rtot}] across
  a transition near $\Mtot^{\star}\sim10^2\,M_\odot$, Eq.~\eqref{eq:massfit}.
\end{itemize}
These are derived for aligned spins with the \texttt{IMRPhenomD} waveform and a
single LIGO design-sensitivity detector; the framework is configurable for
arbitrary detector networks. Precessing spins (and the corresponding
generalisation of $\chi_p$) require a precessing waveform model and are left to
future work.

\end{document}
